\section{Domain Task 10: Trải nghiệm khách hàng thông qua Khai phá văn bản}

Câu hỏi nghiệp vụ: Khách hàng thường đề cập đến những từ khóa nào nhất?
Những yếu tố cốt lõi nào tạo nên sự hài lòng (Điểm cao) và yếu tố nào dẫn
đến sự phàn nàn (Điểm thấp)?

Mục đích: Ứng dụng Khai phá văn bản (NLP) để lượng hóa cảm xúc từ bình
luận, giúp chủ nhà / nền tảng biết chính xác các điểm cần tối ưu.

\subsection{Biểu đồ 2: Yếu tố thúc đẩy sự hài lòng và phàn nàn (Diverging Bar Chart)}

\subsubsection*{A. Idiom}

\begin{longtable}{@{} p{2.8cm} p{12cm} @{}}
\toprule
\textbf{Đặc điểm} & \textbf{Chi tiết} \\
\midrule
\endhead

\bottomrule
\endfoot

\textbf{Idiom} & Diverging Bar Chart (Biểu đồ thanh phân kỳ) \\
\addlinespace[0.2cm]

\textbf{What} & Từ khóa thuộc Whitelist (Keyword): C, Nhóm đánh giá (Dominant Group): C, Độ lệch đặc trưng (Strength / Log-Odds Ratio): Q \\
\addlinespace[0.2cm]

\textbf{Why} & produce (derive) $\rightarrow$ locate $\rightarrow$ compare
\begin{itemize}[nosep, leftmargin=*, topsep=0.1cm]
    \item Derive: Tính toán chỉ số toán học Strength (Log-Odds Ratio) để đo lường mức độ đặc trưng của cụm từ thuộc về nhóm Điểm Cao hay Điểm Thấp (loại bỏ nhiễu do mất cân bằng dữ liệu lượng review).
    \item Xác định (locate) chính xác các yếu tố khiến khách ghét và yếu tố khiến khách khen để so sánh (compare) mức độ nghiêm trọng.
\end{itemize} \\
\addlinespace[0.2cm]

\textbf{How} & \textbf{Encode:}
\begin{itemize}[nosep, leftmargin=*, topsep=0.1cm]
    \item Mark: Line (bar mark)
    \item Channel:
    \begin{itemize}[nosep, leftmargin=0.5cm]
        \item PosX: Trục X phân kỳ ở giữa (0). Đâm sang phải (Dương) thể hiện sự Hài lòng. Đâm sang trái (Âm) thể hiện sự Phàn nàn.
        \item PosY: Danh sách từ khóa.
        \item Color: Phân biệt Nhóm Đánh giá (Xanh = Hài lòng, Cam = Phàn nàn).
    \end{itemize}
\end{itemize}

\vspace{0.2cm}
\textbf{Reduce:}
\begin{itemize}[nosep, leftmargin=*, topsep=0.1cm]
    \item Filter: Áp dụng Whitelist để chỉ giữ lại các n-grams thực sự mang ý nghĩa Insight (loại bỏ tên riêng, địa danh).
    \item Sort: Sắp xếp giảm dần theo chỉ số Strength trên trục Y.
\end{itemize} \\
\end{longtable}

\begin{figure}[H]
    \centering
    \safeincludegraphics[width=0.9\linewidth]{charts/domain_task10_chart2.png}
    \caption{Hình 10.2. Biểu đồ yếu tố thúc đẩy sự hài lòng và phàn nàn}
    \label{fig:domain-task10-chart2}
\end{figure}

\subsubsection*{B. Đánh giá biểu đồ}

\textbf{1. Nguyên lý biểu đạt:}

\begin{itemize}
    \item Thuộc tính: Keyword (C), Nhóm (C), Strength (Q).
    \item Channel:
    \begin{itemize}
        \item PosX: Đâm về 2 hướng $\rightarrow$ Biểu đạt tính chất đối lập (Dương/Âm) của chỉ số Log-Odds.
        \item PosY: Danh sách từ khóa.
        \item Color: Xanh/Cam $\rightarrow$ Mang tính ngữ nghĩa cảm xúc rất chuẩn xác.
    \end{itemize}
    \item Đánh giá mức độ quan trọng:
    \begin{itemize}
        \item Ưu tiên 1 -- Chiều dài và Hướng (PosX / Length): Trọng tâm của biểu đồ là độ lệch Strength (Q). Theo Mackinlay, việc gán (Q) vào kênh Position/Length trên thang đo chung là lựa chọn số 1 về độ chính xác. Trục giữa (0) phân kỳ ra 2 bên làm tăng khả năng biểu đạt tính đối lập.
        \item Ưu tiên 2 -- Màu sắc (Hue Color): Thuộc tính nhóm đánh giá (Hài lòng/Phàn nàn -- C) được gán cho kênh Hue. Phù hợp hoàn hảo với định luật Mackinlay cho dữ liệu Categorical, đồng thời mang ý nghĩa cảnh báo về mặt cảm xúc.
        \item Ưu tiên 3 -- Vị trí dọc (PosY): Dùng để xếp hạng các từ khóa Keyword (C).
    \end{itemize}
\end{itemize}

\textbf{2. Nguyên lý hiệu quả:}

\begin{itemize}
    \item Độ chính xác (Accuracy): Việc áp dụng thuật toán NLP (Log-Odds cơ số 2) thay vì đếm Count thông thường đã giải quyết triệt để sự chênh lệch giữa lượng review tốt và xấu. Chiều dài cột (Length) tuân thủ định luật Stevens ($n = 1.0$), cung cấp độ chính xác, Log\_error vô cùng thấp.
    \item Khả năng phân biệt (Discriminability): Cực kỳ cao. Trục Diverging chẻ làm 2 phía trái/phải kết hợp màu Xanh (an toàn) và Cam (cảnh báo) giúp não bộ người xem phân loại rạch ròi vùng Khen và vùng Chê ngay lập tức mà không gặp trở ngại nào.
    \item Khả năng tách biệt (Separability): Kênh PosX, PosY và Hue Color kết hợp đồng nhất, tách biệt hoàn hảo, tối ưu hóa quá trình nhận thức thị giác.
\end{itemize}

\subsubsection*{C. Phân tích biểu đồ (Insight)}

\begin{itemize}
    \item Yếu tố tạo sự Hài lòng (màu Xanh): Khách hàng bị chinh phục hoàn toàn bởi tính thẩm mỹ và vệ sinh tuyệt đối (các từ khóa đứng đầu là ``beautifully decorated'', ``immaculate'', ``gorgeous''). Các dịch vụ gia tăng như ``stocked kitchen'' hay ``warm welcome'' là chìa khóa giúp chủ nhà đạt điểm tuyệt đối.
    \item Yếu tố gây Phàn nàn (màu Cam): Biểu đồ phơi bày các yếu tố khiến điểm số sụt giảm. Tiêu biểu nhất là các vấn đề chi phí ẩn (``hidden fees'', ``security deposit''). Kế đến là thái độ dịch vụ (``terrible experience'', ``unresponsive'', ``scam'', ``refused'') và tệ nhất là vấn đề môi trường phòng (``infestation'', ``smelly'', ``cigarette''). Khuyến nghị chủ nhà cần minh bạch về giá cả và khử mùi/côn trùng lập tức.
\end{itemize}
